\documentclass[11pt,a4paper]{article}
% =====================================================================
%  Shared preamble for the Part V lecture notes (lectures 19-22).
%
%  These four lectures sit outside Geron, so the notes ARE the primary
%  source and are examined on the same terms as the textbook parts.
%  Everything here is chosen to make that readable on paper: one column,
%  generous measure, and the same six-colour palette as the slides so a
%  student moving between the two is not re-learning what red means.
%
%  Build with pdflatex (twice, for the toc and the refs):
%      cd notes && latexmk -pdf lecture-19.tex
%  or  make -C notes
% =====================================================================
\usepackage[T1]{fontenc}
\usepackage[utf8]{inputenc}
\usepackage{newtxtext,newtxmath}      % Times-like: prints well, reads well
\usepackage[scaled=0.86]{inconsolata} % code, at a size that sits beside Times
\usepackage{microtype}
% newtx defines \Bbbk, \openbox, \proof and \endproof; amssymb and amsthm
% then redefine them and abort. Freeing the four names before those two
% packages load is the standard fix, and it costs nothing here: we use
% amsthm only for \newtheorem.
\let\Bbbk\relax
\usepackage{amsmath,amssymb}
\let\openbox\relax
\let\proof\relax
\let\endproof\relax
\usepackage{amsthm}
\usepackage{booktabs}
\usepackage{enumitem}
\usepackage{xcolor}
\usepackage{tcolorbox}
\usepackage{titlesec}
\usepackage{graphicx}
\usepackage[hidelinks]{hyperref}
\tcbuselibrary{skins,breakable}

% ---- the course palette, from assets/css/custom.css -------------------
\definecolor{aimlprimary}{HTML}{0B3D62}   % deep blue  - structure
\definecolor{aimlaccent} {HTML}{C0392B}   % brick red  - failures
\definecolor{aimlsuccess}{HTML}{14663A}   % green      - repairs
\definecolor{aimlmath}   {HTML}{6C3483}   % purple     - the derivation
\definecolor{aimlmuted}  {HTML}{4B5563}
\definecolor{aimlrule}   {HTML}{B0BCC7}
\definecolor{aimlsoft}   {HTML}{F4F7F9}

\usepackage[a4paper,margin=32mm,top=30mm,bottom=30mm]{geometry}
\setlength{\parskip}{0.45\baselineskip}
\setlength{\parindent}{0pt}
\linespread{1.06}

% ---- headings --------------------------------------------------------
\titleformat{\section}{\Large\bfseries\color{aimlprimary}}{\thesection}{0.7em}{}
\titleformat{\subsection}{\large\bfseries\color{aimlprimary}}{\thesubsection}{0.7em}{}
\titleformat{\subsubsection}{\bfseries\color{aimlmuted}}{}{0em}{}
\titlespacing*{\section}{0pt}{1.6\baselineskip}{0.5\baselineskip}
\titlespacing*{\subsection}{0pt}{1.2\baselineskip}{0.35\baselineskip}

% ---- boxes -----------------------------------------------------------
% One box per role, and the role is the colour. A student who has seen the
% slides already knows what each one means before reading a word of it.
\newtcolorbox{keybox}[1][]{
  breakable, enhanced, colback=aimlsoft, colframe=aimlprimary,
  boxrule=0.4pt, left=8pt, right=8pt, top=6pt, bottom=6pt,
  fonttitle=\bfseries\small, coltitle=white, colbacktitle=aimlprimary,
  attach boxed title to top left={yshift=-2mm, xshift=4mm},
  boxed title style={boxrule=0pt, arc=1pt}, #1}

\newtcolorbox{failbox}[1][]{
  breakable, enhanced, colback=white, colframe=aimlaccent,
  boxrule=0.9pt, left=8pt, right=8pt, top=6pt, bottom=6pt,
  fonttitle=\bfseries\small, coltitle=white, colbacktitle=aimlaccent,
  attach boxed title to top left={yshift=-2mm, xshift=4mm},
  boxed title style={boxrule=0pt, arc=1pt}, #1}

\newtcolorbox{derivbox}[1][]{
  breakable, enhanced, colback=white, colframe=aimlmath,
  boxrule=0.6pt, left=8pt, right=8pt, top=6pt, bottom=6pt,
  fonttitle=\bfseries\small, coltitle=white, colbacktitle=aimlmath,
  attach boxed title to top left={yshift=-2mm, xshift=4mm},
  boxed title style={boxrule=0pt, arc=1pt}, #1}

% An exercise. Deliberately the same shape as the notebook's `try` field:
% one change, and what should happen to the output.
\newtcolorbox{trybox}[1][]{
  breakable, enhanced, colback=white, colframe=aimlsuccess,
  boxrule=0.5pt, left=8pt, right=8pt, top=6pt, bottom=6pt,
  fonttitle=\bfseries\small, coltitle=white, colbacktitle=aimlsuccess,
  attach boxed title to top left={yshift=-2mm, xshift=4mm},
  boxed title style={boxrule=0pt, arc=1pt}, #1}

% ---- small semantics -------------------------------------------------
\newcommand{\scope}[1]{\textcolor{aimlmuted}{\small #1}}
\newcommand{\term}[1]{\textbf{#1}}
\newcommand{\code}[1]{\texttt{\small #1}}
\newcommand{\lecture}[1]{Lecture~#1}

\newtheorem{definition}{Definition}[section]
\newtheorem{proposition}[definition]{Proposition}

% ---- title block -----------------------------------------------------
% \notestitle{19}{Information retrieval: the lexical foundation}{SciFact (BEIR)}
\newcommand{\notestitle}[3]{%
  \thispagestyle{empty}
  \begin{flushleft}
    {\color{aimlmuted}\small\sffamily
     APPLICATIONS OF MACHINE LEARNING \quad$\cdot$\quad
     BSc Mathematics of Artificial Intelligence}\\[2mm]
    {\color{aimlrule}\rule{\textwidth}{0.8pt}}\\[4mm]
    {\color{aimlmuted}\large Lecture #1 \quad$\cdot$\quad Part V \quad$\cdot$\quad
     Extended notes}\\[3mm]
    {\Huge\bfseries\color{aimlprimary}\baselineskip=1.15\baselineskip #2\par}\vspace{4mm}
    {\color{aimlmuted}\large Dataset: #3}\\[3mm]
    {\color{aimlrule}\rule{\textwidth}{0.8pt}}
  \end{flushleft}
  \vspace{2mm}
  \begin{keybox}[title={Why these notes exist}]
    Lectures 19--22 sit outside G\'eron, the course's primary text. For those
    four lectures \emph{these notes are the primary source}, and they are
    examined on exactly the same terms as the textbook parts. Everything in
    the slide deck for this lecture is here, with the arguments written out at
    length rather than compressed onto a slide; the numbers are the ones the
    accompanying notebook prints, and every one of them is a measurement on
    the corpus named above rather than a figure quoted from a paper.
  \end{keybox}
  \vspace{1mm}
  {\small\tableofcontents}
  \vspace{2mm}
  \noindent{\color{aimlrule}\rule{\textwidth}{0.4pt}}
  \vspace{1mm}
}


\begin{document}
\notestitlefor{14}{III}{Detection and segmentation}{COCO, 128 images}{Chapter 12}

\section{When a label is not an answer}

Three lectures of vision, all answering the same question: what is in this
image? Neither a convolutional network nor a pretrained backbone can express
\emph{where} the object is, or \emph{how many} there are. A photograph with
three cyclists has one label and three answers.

A retail shelf audit puts it plainly: \emph{``We photograph shelves. We need to
know how many of each product are on each shelf, and where each one is, so that
a picking route can be planned. A list of the categories present is no use to
us.''}

\begin{center}
\small
\begin{tabular}{lll}
\toprule
& \textbf{Classification} & \textbf{Detection} \\
\midrule
Question & what is in this picture & what is where, and how many \\
Output & one vector, length $K$ & three arrays, length $N$ \\
$N$ is & fixed before you look & decided by the image and by two cut-offs \\
Wanted from pooling & invariance & equivariance \\
\bottomrule
\end{tabular}
\end{center}

The last row is Lecture 12 with a consequence: \emph{a network that has learned
not to care where the object is cannot tell you where it is.} So the output
stops being a label and becomes a set of boxes --- which means the first thing
we need is a new notion of what a correct answer is.

\begin{failbox}[title={A box is four numbers, and which four is not agreed}]
\code{[x1, y1, x2, y2]} is two opposite corners, used by torchvision.
\code{[x, y, w, h]} is a corner and a size, used by the COCO annotation file.
\code{[cx, cy, w, h]} is a centre and a size, used by the YOLO family.

You will hold the first two in the same notebook. Read \code{[100, 50, 20, 30]}
as corners and you get a box 20 wide when it should be 20 \emph{starting at}
100 --- a valid box, in the wrong place, with no error.
\end{failbox}

\subsection{Our corpus is 128 images}

Not COCO. 128 images out of the 5{,}000 in \code{val2017}, chosen as the
numerically lowest image ids --- a rule, not a selection. Full COCO with its
training split is roughly 20 GB and is not downloaded anywhere in this course.

\begin{center}
\small
\begin{tabular}{lr}
\toprule
\textbf{Quantity} & \textbf{Value} \\
\midrule
Images & 128 \\
Annotated objects, crowd regions excluded & 898 \\
Crowd regions, dropped & 8 \\
Distinct categories appearing & 73 of 80 \\
Objects per image --- mean & 7.02 \\
Objects per image --- median & 5 \\
Objects per image --- range & 1 to 29 \\
\bottomrule
\end{tabular}
\end{center}

Mean 7.02 against median 5: a long right tail, and one image alone carries 29
objects. And 39.0\% of every annotated object in this corpus is a
\emph{person} --- remember that when we start averaging over categories.

\emph{Every number in this lecture is measured on 128 images.} When we quote a
score, we quote the sample size beside it.

\section{Derivation: scoring a box, scoring a detector}

\subsection{One quantity, four requirements}

Before writing anything down, fix what it has to do: be 1 for coincident boxes
and 0 for disjoint ones; punish a box for being too large; punish a box for
being too small; and be dimensionless, so a cup and a sofa are on one scale.

Requirement 2 says the predicted area must be in the denominator, requirement 3
says the true area must be too, and requirement 4 says the numerator must be an
area as well. There is essentially one candidate:
\[ \mathrm{IoU}(A,B) = \frac{|A \cap B|}{|A \cup B|}
   = \frac{|A \cap B|}{|A| + |B| - |A \cap B|}. \]
The second form is the one you compute --- the union is never built --- and
subtracting the intersection is not an optimisation: add the two areas and you
have counted the overlap twice.

A box covering the whole image scores 0.086 on our corpus: always right about
something and never about anything in particular.

\begin{keybox}[title={IoU does three different jobs}]
\emph{Matching}: does a detection count as the same object as an annotation.
\emph{Suppression}: are two detections duplicates of each other.
\emph{Training}: as the loss a regression head descends.

The first two use IoU as a comparison. Only the third needs it to have a
derivative, and that is where the derivation has its teeth.
\end{keybox}

\subsection{IoU as a loss, and why it fails}

Two boxes, each 100 by 100. Fix the first at the origin, slide the second right
by $d$ pixels, and record the value and the gradient. Synthetic on purpose ---
nothing here depends on a dataset, so the conclusion is a property of the
formula rather than of COCO.

\begin{center}
\small
\begin{tabular}{rrrrr}
\toprule
$d$ & \textbf{IoU} & $\partial\,\mathrm{IoU}/\partial d$ &
\textbf{GIoU} & $\partial\,\mathrm{GIoU}/\partial d$ \\
\midrule
  0 & 1.000 &  0.000 &  1.000 &  0.000 \\
 50 & 0.333 & $-0.0089$ & 0.333 & $-0.0089$ \\
 90 & 0.053 & $-0.0055$ & 0.053 & $-0.0055$ \\
150 & 0.000 &  0.000 & $-0.200$ & $-0.0032$ \\
300 & 0.000 &  0.000 & $-0.500$ & $-0.0013$ \\
\bottomrule
\end{tabular}
\end{center}

Rows four and five are the point: two boxes 150 pixels apart and two boxes 300
pixels apart are, to IoU, \emph{exactly equally wrong}.

\begin{derivbox}[title={Why the gradient is exactly zero, not merely small}]
For $d \ge 100$ the overlap width is $\max(0, 100-d) = 0$, so
\[ |A \cap B| \equiv 0
   \implies \mathrm{IoU} \equiv 0
   \implies \frac{\partial\,\mathrm{IoU}}{\partial d} \equiv 0. \]
The clamp makes the function \emph{constant} on the whole disjoint region, not
merely small there --- and a constant has no descent direction. Not a weak one:
none.
\end{derivbox}

This is a property of the formula, and no optimiser, learning rate or
initialisation repairs it. A proposal that starts anywhere outside the object
it should find receives the same loss and the same zero gradient wherever it
is, because nothing in the loss knows which way the object lies. Which is why
detectors are not trained on IoU alone, and why the coordinate regression they
\emph{are} trained on needs anchors close enough to overlap in the first place.

The repair is to put the empty space back in. With $C$ the smallest box
containing both,
\[ \mathrm{GIoU}(A,B) = \mathrm{IoU}(A,B)
   - \frac{|C| - |A \cup B|}{|C|}. \]
When the boxes are disjoint and move apart, $|C|$ grows and $|A \cup B|$ does
not, so the penalty grows. CIoU adds two more terms --- a centre distance
normalised by the enclosing diagonal, and an aspect-ratio disagreement --- and
converges faster because it moves the centre directly rather than shrinking a
gap.

\subsection{From one box to a whole detector}

IoU scores one pair. A detector emits 4{,}419 boxes over 128 images, most of
them wrong, all carrying a score. Fixing a score threshold gives one precision
and one recall --- and throws away everything the ranking knew.

\emph{A detector is a ranking. Score it as one.} Fix a class and an IoU
threshold, sort every detection of that class over the whole corpus by score,
and walk down the list: for each detection find the best \emph{unmatched}
annotation in the same image, and if that IoU clears the threshold it is a true
positive and the annotation is now matched. That word \emph{unmatched} is what
makes a second box on the same bottle a false positive rather than a second
success.

Precision and recall are unchanged from Lecture 3, computed cumulatively after
each detection --- so a corpus of 1{,}209 person detections gives 1{,}209
points. The recall denominator is fixed, so recall is non-decreasing; precision
has both a numerator and a denominator moving.

\begin{keybox}[title={Every step accounted for, exactly}]
Walking the person ranking one detection at a time:

\begin{center}
\small
\begin{tabular}{lr}
\toprule
& \textbf{Steps} \\
\midrule
the detection is a false positive --- precision falls & 899 \\
a true positive and precision was below 1 --- precision rises & 254 \\
a true positive and precision was already 1 --- no change & 55 \\
\bottomrule
\end{tabular}
\end{center}

899 is exactly the number of false positives, and $254 + 55 = 309$ is the 310
true positives less the one at the top of the ranking, which has no step above
it. The 55 flat steps are one contiguous run: the top 56 person detections in
the whole corpus are all correct.

An exact classification of all 1{,}208 steps, in the same form as Lecture 3's.
\end{keybox}

Lecture 3 promised the repair: average precision is defined using the maximum
precision at or above each recall level,
\[ p_{\text{env}}(r) = \max_{\tilde r \ge r} p(\tilde r), \]
non-increasing by construction. Why a maximum rather than a smoothing? It has
an operational reading --- if you will accept recall $r$, here is the best
precision available at that recall or better; it is monotone by construction,
so the area under it does not depend on where you sampled; and it is one line
of code with no parameters.

\[ \mathrm{AP} = \int_0^1 p_{\text{env}}(r)\,\mathrm{d}r
   = \sum_k (r_k - r_{k-1})\,p_{\text{env}}(r_k). \]
The sum only has terms where recall actually changed, which is exactly at the
true positives --- and \emph{no threshold appears anywhere}. That is the whole
point: the number does not depend on a cut-off nobody chose.

\begin{center}
\small
\begin{tabular}{lr}
\toprule
\textbf{person, IoU at least 0.5} & \textbf{Value} \\
\midrule
Annotated instances & 350 \\
Detections in the ranking & 1{,}209 \\
True positives & 310 \\
Highest recall reached & 0.886 \\
AP & 0.751 \\
\bottomrule
\end{tabular}
\end{center}

Recall stops at 0.886 because 40 annotated people are never detected at all, at
any score. No threshold recovers them, and AP is honest about that.

\subsection{A mean of a mean of an area}

The first mean is over classes: $\mathrm{mAP}(t)$ averages $\mathrm{AP}_c(t)$
over the classes present. Unweighted --- a class with one annotation counts as
much as person with 350. We average over the 73 classes appearing in our 128
images; the other 7 have no annotations here. \emph{Which classes are in the
mean is a decision, it changes the number, and papers do not always say which
one they made.} Thirteen classes score exactly 1.000, and they are the ones
with one or two instances in the whole corpus.

The second mean is over IoU thresholds, averaging the ten values from 0.50 to
0.95. At 0.50 a box that half-covers the object counts as correct; at 0.95 it
must be very nearly exact. That average is COCO's headline number and what
``mAP'' means unqualified in a modern paper --- and on our corpus AP runs from
0.659 at the loosest threshold to 0.040 at the tightest, with one number in the
middle standing for both.

\begin{center}
\small
\begin{tabular}{lr}
\toprule
\textbf{Our detector, on 128 images} & \textbf{Value} \\
\midrule
mAP at IoU 0.50 & 0.659 \\
mAP at IoU 0.75 & 0.475 \\
mAP averaged over 0.50 to 0.95 & 0.439 \\
\bottomrule
\end{tabular}
\end{center}

Three collapses, one number. Each throws away a different thing, and none is
recoverable from the result.

\section{Running a detector, and reading it}

The specification first: 128 JPEG files and one annotation file in; per image,
boxes in corner form, category ids, scores, and one integer count out; nothing
trained; and the \emph{check} --- every predicted box lies inside the image,
scores are non-increasing, and the count equals the number of boxes kept. That
last row is the one people skip, and it is the one that catches a
corner-versus-size mix-up in a line.

\begin{keybox}[title={Convert once, at the edge, and assert it}]
COCO gives $x, y, w, h$; everything else here is corners. The conversion
happens once, on load, and then:
\begin{verbatim}
assert (b[:, 2] >= b[:, 0]).all(), "x2 < x1 — you read w as x2"
assert (b[:, 3] >= b[:, 1]).all(), "y2 < y1 — same bug"
\end{verbatim}
Reading $w$ as $x_2$ gives $x_2 < x_1$ for every box wider than its own left
edge, which is almost all of them. Two lines catch it; nothing else will.
\end{keybox}

One photograph returns 88 boxes --- not 88 objects. The three arrays share an
index, and they arrive sorted by score descending, which is a thing to assert
rather than assume.

\begin{failbox}[title={There is already a cut-off, and it is not yours}]
\code{box\_score\_thresh} defaults to 0.05, so anything below is never returned
at all. \code{box\_nms\_thresh} defaults to 0.5, so near-duplicates are removed
before you see them. \code{box\_detections\_per\_img} defaults to 100 --- and
our busiest image returns exactly 100 boxes, so that cap is not hypothetical.

Reviewer question 5, three times over.
\end{failbox}

The count is a function of the threshold: move the cut-off and the answer to
the stakeholder's question changes by a factor of ten. We adopt 0.5, because it
is the conventional reporting threshold \emph{and because it was chosen before
looking at the error curve}. The minimum of that curve sits at 0.75 with an MAE
of 1.92, and adopting it would be choosing a hyperparameter on the data we then
report --- Lecture 2's error in a new costume. So it is quoted as a diagnostic
and not as a result.

\begin{center}
\small
\begin{tabular}{lr}
\toprule
\textbf{At score at least 0.5} & \textbf{Value} \\
\midrule
Count MAE & 3.00 \\
Signed count error & $+2.41$ \\
Images counted exactly right & 27 of 128 \\
Images with too many boxes & 83 \\
Images with too few & 18 \\
\bottomrule
\end{tabular}
\end{center}

Two numbers, not one: the absolute error says how far off, the signed error
says which way --- and the cloud sits above the diagonal, so the failure mode
is too many boxes. 3.00 against a baseline of 6.02, so the model halves the
error of a system that never looks at the image.

\subsection{Check your own implementation against someone else's}

We derived AP here and implemented it ourselves, which is exactly the situation
in which a result should not be believed.

\begin{center}
\small
\begin{tabular}{lrr}
\toprule
\textbf{mAP on our 128 images} & \textbf{Ours} & \textbf{Reference evaluator} \\
\midrule
at IoU 0.50 & 0.659 & 0.660 \\
at IoU 0.75 & 0.475 & 0.476 \\
averaged 0.50 to 0.95 & 0.439 & 0.440 \\
\bottomrule
\end{tabular}
\end{center}

Largest disagreement 0.0014, against the official COCO evaluator run on the
same predictions by a library we did not write.

\begin{failbox}[title={And now compare 128 images against 5{,}000}]
Our averaged mAP is 43.9 on 128 images. torchvision reports 37.0 for
\emph{these same weights} on all 5{,}000.

6.9 points higher, and nothing about the detector changed. The difference is
the sample: 128 images, many of them easy, and 73 classes several of which have
a single instance that happened to be found. This is what a small evaluation
set does --- and it does it in the flattering direction more often than not.
\end{failbox}

\subsection{Suppression}

Without suppression the detector proposes 231.7 candidate boxes per image,
against 23.6 after suppression at IoU 0.5, and 7.02 actual objects. \emph{The
detector you ran had already thrown away nine tenths of its own output before
you saw it, using IoU, at a threshold you did not set.}

Use \code{batched\_nms} rather than \code{nms}: a bottle overlapping a glass
must not suppress it, because they are different classes and both are true. And
the default of 0.5 happens to be near the best value on our corpus --- a fact
about this corpus, not a law.

\section{Beyond boxes}

A box was always an approximation. A bicycle's box is mostly not bicycle; a
person sitting cross-legged fills perhaps half of theirs; and for a picking arm
``somewhere in this rectangle'' may not be good enough.

\term{Semantic segmentation} gives one class label per pixel, so two people
standing together are one person region and you cannot count them. The
architectural problem is that pooling destroyed the resolution the answer needs
--- Lecture 12's invariance arriving as a cost --- and the repair is upsampling
with skip connections that recover detail from the earlier, higher-resolution
layers.

\term{Instance segmentation} gives one binary mask per detected object. Mask
R-CNN is Faster R-CNN with one extra head predicting a small mask inside each
proposed box, and it is two lines away from what you already ran. Only that
output can answer ``how many''.

Scoring a mask is the same argument again: mask AP replaces box IoU with mask
IoU, intersecting pixel sets instead of rectangles, and everything above ---
the ranking, the envelope, the area, the two means --- is unchanged.
torchvision reports 34.6 mask mAP for these weights against 37.9 box mAP on the
full 5{,}000, which is what ``a mask is harder than a box'' means numerically.

And detection on video is not tracking. Run the detector on every frame and you
have boxes, not the same bottle in frame 2 as in frame 1. Tracking asks for a
persistent identity, which turns each frame into an assignment problem --- and
the cheapest similarity between a track's predicted box and a new detection is
IoU, so the derivation pays for itself a third time. It also fails the same
way: an object occluded for ten frames reappears where its predicted box does
not overlap, IoU is zero, the assignment finds nothing, and the object is
issued a new identity.

\section{The notebook}

\begin{trybox}[title={What to change}]
Move one predicted box further and further from its match, and watch IoU sit at
exactly zero the whole way. That flat zero is the vanishing gradient the
derivation predicts --- and seeing it as a constant rather than as a small
number is the difference between the two answers to the exam question.
\end{trybox}

\section{Exercises}

\begin{enumerate}[leftmargin=1.6em,itemsep=4pt]
  \item Two boxes are disjoint and moved further apart. State what IoU and its
        derivative do, and why no learning rate repairs it. \hfill \scope{[5]}
  \item An IoU implementation without a clamp is tested on overlapping boxes
        only and passes. Give the input that breaks it, and the property test
        that would have caught it. \hfill \scope{[5]}
  \item Average precision is defined with a maximum inside it. State what that
        maximum repairs, and which earlier lecture proved the problem.
        \hfill \scope{[4]}
  \item A detector is evaluated on a corpus in which 30 of the classes present
        have exactly one annotated instance each, and reports mAP 0.72 at IoU
        0.5. Give two distinct reasons the number may be optimistic, and one
        thing you would report alongside it. \hfill \scope{[3]}
  \item Explain why semantic segmentation cannot answer ``how many'', and name
        the architectural property that causes it. \hfill \scope{[4]}
\end{enumerate}

\section*{Summary}
\addcontentsline{toc}{section}{Summary}

A label is not an answer when the question is \emph{where} and \emph{how many},
and a network trained to be invariant to position cannot supply either. IoU is
the essentially unique quantity meeting the four requirements --- and it is
exactly zero for every disjoint pair, so as a loss it offers no gradient
however far apart the boxes are. Not a small one: none, because the clamp makes
it constant. GIoU and CIoU repair that by adding a term that keeps moving.

Average precision is built by matching a ranking against annotations, and its
definition contains a maximum for one reason: precision is not monotone, which
Lecture 3 proved, and the envelope is the repair. Counted on the person
ranking, 899 steps down and 254 up and 55 flat account for all 1{,}208 exactly.
mAP then averages over classes and over ten IoU thresholds --- a mean of a mean
of an area, where our own AP runs from 0.659 to 0.040 across the thresholds it
averages.

Our implementation agrees with the official evaluator to 0.0014, which is why
it can be quoted. And it reports 43.9 where torchvision reports 37.0 for the
same weights: 6.9 points, entirely from evaluating on 128 images instead of
5{,}000, in the flattering direction.

\end{document}
